\documentclass[10pt,a4paper]{amsart}
\usepackage{amsmath,amssymb,amsthm,graphicx}
% Retained preamble for the Stacks Project chapter {CH} (shared project preamble).
% The chapter file pinned-original-source.tex only carries \input{{preamble}} on its first line,
% so no preamble text lives inside the chapter itself.  The Stacks Project uses one single shared
% preamble for every chapter: the snapshot stages it three times (sources/stacks-chow/preamble.tex,
% sources/stacks-etale-cohomology/preamble.tex and focused-seeds/AGP/sources/stacks-snapshot-20261011/
% preamble.tex) and those copies are identical in size.  The block below is copied mechanically,
% line for line, from sources/stacks-chow/preamble.tex lines 14-18 and 154-217: the Xy-pic setup and
% the native theorem declarations together with the mathematical macros.  Nothing was added, removed
% or redefined; the wrapper class and the amsmath/amssymb/amsthm/graphicx base packages are supplied
% by the assembler.
% For commutative diagrams we use Xy-pic
% Chapter: stacks-spaces-more-morphisms.tex -- More on Morphisms of Spaces.
\usepackage[all]{xy}

% We use 2cell for 2-commutative diagrams.
\xyoption{2cell}

% Theorem environments.
%
\theoremstyle{plain}
\newtheorem{theorem}[subsection]{Theorem}
\newtheorem{proposition}[subsection]{Proposition}
\newtheorem{lemma}[subsection]{Lemma}

\theoremstyle{definition}
\newtheorem{definition}[subsection]{Definition}
\newtheorem{example}[subsection]{Example}
\newtheorem{exercise}[subsection]{Exercise}
\newtheorem{situation}[subsection]{Situation}

\theoremstyle{remark}
\newtheorem{remark}[subsection]{Remark}
\newtheorem{remarks}[subsection]{Remarks}

\numberwithin{equation}{subsection}

% Macros
%
\def\lim{\mathop{\mathrm{lim}}\nolimits}
\def\colim{\mathop{\mathrm{colim}}\nolimits}
\def\Spec{\mathop{\mathrm{Spec}}}
\def\Hom{\mathop{\mathrm{Hom}}\nolimits}
\def\Ext{\mathop{\mathrm{Ext}}\nolimits}
\def\SheafHom{\mathop{\mathcal{H}\!\mathit{om}}\nolimits}
\def\SheafExt{\mathop{\mathcal{E}\!\mathit{xt}}\nolimits}
\def\Sch{\mathit{Sch}}
\def\Mor{\mathop{\mathrm{Mor}}\nolimits}
\def\Ob{\mathop{\mathrm{Ob}}\nolimits}
\def\Sh{\mathop{\mathit{Sh}}\nolimits}
\def\NL{\mathop{N\!L}\nolimits}
\def\CH{\mathop{\mathrm{CH}}\nolimits}
\def\proetale{{pro\text{-}\acute{e}tale}}
\def\etale{{\acute{e}tale}}
\def\QCoh{\mathit{QCoh}}
\def\Ker{\mathop{\mathrm{Ker}}}
\def\Im{\mathop{\mathrm{Im}}}
\def\Coker{\mathop{\mathrm{Coker}}}
\def\Coim{\mathop{\mathrm{Coim}}}

% Boxtimes
%
\DeclareMathSymbol{\boxtimes}{\mathbin}{AMSa}{"02}

%
% Macros for moduli stacks/spaces
%
\def\QCohstack{\mathcal{QC}\!\mathit{oh}}
\def\Cohstack{\mathcal{C}\!\mathit{oh}}
\def\Spacesstack{\mathcal{S}\!\mathit{paces}}
\def\Quotfunctor{\mathrm{Quot}}
\def\Hilbfunctor{\mathrm{Hilb}}
\def\Curvesstack{\mathcal{C}\!\mathit{urves}}
\def\Polarizedstack{\mathcal{P}\!\mathit{olarized}}
\def\Complexesstack{\mathcal{C}\!\mathit{omplexes}}
% \Pic is the operator that assigns to X its picard group, usage \Pic(X)
% \Picardstack_{X/B} denotes the Picard stack of X over B
% \Picardfunctor_{X/B} denotes the Picard functor of X over B
\def\Pic{\mathop{\mathrm{Pic}}\nolimits}
\def\Picardstack{\mathcal{P}\!\mathit{ic}}
\def\Picardfunctor{\mathrm{Pic}}
\def\Deformationcategory{\mathcal{D}\!\mathit{ef}}
\newenvironment{reference}{}{}
\newenvironment{slogan}{}{}
\newenvironment{history}{}{}
\usepackage{cleveref}
\crefname{originallemma}{original lemma}{original lemmas}
\Crefname{originallemma}{Original lemma}{Original lemmas}
\makeatletter
\newlabel{coherent-lemma-inverse-systems-abelian}{{087X}{source}{}{sourcepointer-1}{}}
\newlabel{coherent-lemma-inverse-systems-abelian@cref}{{[originallemma][0][]087X}{[1][0][]source}{}{}{}}
\makeatother
\crefname{originallemma}{original lemma}{original lemmas}
\Crefname{originallemma}{Original lemma}{Original lemmas}
\makeatletter
\newlabel{coherent-lemma-inverse-systems-pullback}{{0887}{source}{}{sourcepointer-2}{}}
\newlabel{coherent-lemma-inverse-systems-pullback@cref}{{[originallemma][0][]0887}{[1][0][]source}{}{}{}}
\makeatother
\crefname{originallemma}{original lemma}{original lemmas}
\Crefname{originallemma}{Original lemma}{Original lemmas}
\makeatletter
\newlabel{groupoids-lemma-abelian}{{077S}{source}{}{sourcepointer-3}{}}
\newlabel{groupoids-lemma-abelian@cref}{{[originallemma][0][]077S}{[1][0][]source}{}{}{}}
\makeatother
\crefname{originallemma}{original lemma}{original lemmas}
\Crefname{originallemma}{Original lemma}{Original lemmas}
\makeatletter
\newlabel{spaces-morphisms-lemma-flat-local}{{03MM}{source}{}{sourcepointer-4}{}}
\newlabel{spaces-morphisms-lemma-flat-local@cref}{{[originallemma][0][]03MM}{[1][0][]source}{}{}{}}
\makeatother
\title{The category of coherent sheaves with an I-adic filtration structure on a Noetherian algebraic space is abelian, with exactness etale-local and flat-pullback exact}
\author{The Stacks Project Authors}
\date{stacks-project:spaces-more-morphisms.tex:local-snapshot-20261011; GNU FDL 1.2 or later}
\begin{document}
\maketitle
\noindent\textit{Editorial note.} Statement span is lines 9562-9574 of sources/stacks-spaces-more-morphisms/pinned-original-source.tex; proof span is lines 9576-9612, the complete author proof (35 body lines).\par\medskip
\noindent\textit{Editorial note.} Navigation stubs: 4 cited-but-not-included labels are supplied as non-load-bearing pointers with their pinned tags where the snapshot tag index has them.\par\medskip
\noindent\textit{Editorial note.} extra\_packages is empty on purpose: the retained preamble already loads xy with options and a second declaration clashes in the assembler wrapper.\par\medskip
\noindent\textit{Editorial note.} Text is retained verbatim, including the author's own typos and placeholders.\par\medskip
\noindent\textit{Editorial note.} The schemes chapter carries a same-named lemma used as input; it lies in Cohomology of Schemes, outside both pools, and its card status (if any) is another lane's.\par\medskip
\noindent\textit{Allowed standard prerequisites (editorial).}\begin{itemize}
\item Explicit prerequisite boundary. Scheme-side theory of inverse systems of coherent modules (Cohomology of Schemes, lemmas inverse-systems-abelian, inverse-systems-pullback), etale descent for coherent data as in the chapter's preceding discussion, and the groupoid-descent formal template (Groupoids, lemma-abelian).
\end{itemize}\medskip
\ifdefined\hypertarget\hypertarget{sourcepointer-1}{}\fi
\noindent\textit{Source reference.} Original lemma 087X: sources/stacks-spaces-more-morphisms/pinned-original-source.tex, source lines 9612--9612. Cited by the retained statement or proof but not part of the collected spans; supplied only as a navigation pointer. The argument does not depend on its content.\par\medskip
\ifdefined\hypertarget\hypertarget{sourcepointer-2}{}\fi
\noindent\textit{Source reference.} Original lemma 0887: sources/stacks-spaces-more-morphisms/pinned-original-source.tex, source lines 9612--9612. Cited by the retained statement or proof but not part of the collected spans; supplied only as a navigation pointer. The argument does not depend on its content.\par\medskip
\ifdefined\hypertarget\hypertarget{sourcepointer-3}{}\fi
\noindent\textit{Source reference.} Original lemma 077S: sources/stacks-spaces-more-morphisms/pinned-original-source.tex, source lines 9612--9612. Cited by the retained statement or proof but not part of the collected spans; supplied only as a navigation pointer. The argument does not depend on its content.\par\medskip
\ifdefined\hypertarget\hypertarget{sourcepointer-4}{}\fi
\noindent\textit{Source reference.} Original lemma 03MM: sources/stacks-spaces-more-morphisms/pinned-original-source.tex, source lines 9612--9612. Cited by the retained statement or proof but not part of the collected spans; supplied only as a navigation pointer. The argument does not depend on its content.\par\medskip
\begin{lemma}
\label{lemma-inverse-systems-abelian}
Let $S$ be a scheme. Let $X$ be a Noetherian algebraic space over $S$ and
let $\mathcal{I} \subset \mathcal{O}_X$ be a quasi-coherent sheaf of ideals.
\begin{enumerate}
\item The category $\textit{Coh}(X, \mathcal{I})$ is abelian.
\item Exactness in $\textit{Coh}(X, \mathcal{I})$
can be checked \'etale locally.
\item For any flat morphism $f : X' \to X$ of Noetherian algebraic spaces
the functor $f^* : \textit{Coh}(X, \mathcal{I}) \to
\textit{Coh}(X', f^{-1}\mathcal{I}\mathcal{O}_{X'})$ is exact.
\end{enumerate}
\end{lemma}

\begin{proof}
Proof of (1). Choose an affine scheme $U_0$ and a surjective \'etale morphism
$U_0 \to X$. Set $U_1 = U_0 \times_X U_0$ and
$U_2 = U_0 \times_X U_0 \times_X U_0$ as in our discussion of
\'etale descent above. The categories $\textit{Coh}(U_i, \mathcal{I}_i)$
are abelian
(Cohomology of Schemes, Lemma \ref{coherent-lemma-inverse-systems-abelian})
and the pullback functors are exact functors
$\textit{Coh}(U_0, \mathcal{I}_0) \to \textit{Coh}(U_1, \mathcal{I}_1)$
and
$\textit{Coh}(U_1, \mathcal{I}_1) \to \textit{Coh}(U_2, \mathcal{I}_2)$
(Cohomology of Schemes, Lemma \ref{coherent-lemma-inverse-systems-pullback}).
The lemma then follows formally from the description of
$\textit{Coh}(X, \mathcal{I})$ as a category of descent data.
Some details omitted; compare with the proof of
Groupoids, Lemma \ref{groupoids-lemma-abelian}.

\medskip\noindent
Part (2) follows immediately from the discussion in the previous paragraph.
In the situation of (3) choose a commutative diagram
$$
\xymatrix{
U' \ar[d] \ar[r] & U \ar[d] \\
X' \ar[r] & X
}
$$
where $U'$ and $U$ are affine schemes and the vertical morphisms are
surjective \'etale. Then $U' \to U$ is a flat morphism of Noetherian
schemes (Morphisms of Spaces, Lemma \ref{spaces-morphisms-lemma-flat-local})
whence the pullback functor
$\textit{Coh}(U, \mathcal{I}\mathcal{O}_U) \to
\textit{Coh}(U', \mathcal{I}\mathcal{O}_{U'})$
is exact by
Cohomology of Schemes, Lemma \ref{coherent-lemma-inverse-systems-pullback}.
Since we can check exactness in $\textit{Coh}(X, \mathcal{O}_X)$
on $U$ and similarly for $X', U'$ the assertion follows.
\end{proof}

\end{document}
